\section{Literature Review}\label{sec:literature}

Traffic monitoring systems combine sensing and prediction. This section reviews
the two lines of work that inform the prototype: vehicle detection in traffic
imagery and spatio-temporal traffic forecasting. It also covers the evaluation
practices used to compare the resulting models.

\subsection{Vehicle Detection in Traffic Imagery}\label{sec:literature-detection}

One-stage detectors are the usual choice for traffic imagery because a single
forward pass returns boxes and classes at near-real-time speed
\parencite{ultralytics2026yolo}. COCO-pretrained weights are a strong starting
point \parencite{lin2014coco}, but the COCO camera geometry differs from fixed,
high-mounted traffic cameras. MIO-TCD provides a traffic-specific benchmark with
137,743 human-labeled localization frames from thousands of traffic cameras
\parencite{luo2018miotcd}. When target-domain
labels are unavailable, pseudo-labeling is a common approach: a teacher model
labels unlabeled images, and a student is trained on the resulting hard labels
\parencite{lee2013pseudo}. When the same teacher also supplies the evaluation
reference, the resulting score measures agreement with the teacher; the LTA
results in this report are interpreted on that basis.

\subsection{Traffic Forecasting}\label{sec:literature-forecasting}

Spatio-temporal graph models are widely studied for road-network forecasting.
STGCN combines temporal convolution with graph convolution
\parencite{yu2018stgcn}; Graph WaveNet learns an adaptive adjacency and uses
dilated temporal convolutions \parencite{wu2019gwnet}; MTGNN learns a graph and a
multi-step curriculum \parencite{wu2020mtgnn}; AGCRN uses node-specific
embeddings instead of a fixed adjacency \parencite{bai2020agcrn}. LibCity
provides a common implementation and evaluation environment
\parencite{wang2021libcity,jiang2023libcity}. In comparative studies, the models
are often trained with different input channels, optimizers and early-stopping
rules, so an observed difference reflects both architecture and training
recipe. The eight models compared in this project span feed-forward, recurrent,
fixed-graph, learned-graph and node-embedding designs, which gives the benchmark
coverage of the main modeling families rather than a single architecture.

\subsection{Evaluation and Benchmarking}\label{sec:literature-evaluation}

Frames from the same camera are highly correlated, so random image-level splits
leak near-duplicates into the test set. We group images by scene before
splitting and resample whole groups for confidence intervals. Detection accuracy
is reported with COCO-style mAP \parencite{lin2014coco}; uncertainty is estimated
with bootstrap resampling \parencite{efron1993bootstrap}. Forecasting comparisons
are most informative when every model receives the same input information; the
benchmark therefore reports both native configurations and a unified
nine-channel setting.

\subsection{Synthesis and Research Gap}\label{sec:literature-gap}

Across these lines of work, detection and forecasting are usually evaluated on
different data representations: annotated images for recognition and calibrated
sensor series for forecasting. Detection studies also face a cross-camera
adaptation problem when the target cameras differ in viewpoint, scale and
lighting from the training benchmark. Forecasting comparisons add a second
problem: native configurations differ in input features and training recipes, so
a ranking mixes architecture with implementation choices. Monitoring
integration raises a third problem: the two components produce different
outputs, at different time scales, and their data domains must remain visible to
the user. These gaps motivate a modular prototype that evaluates each component
on its own domain, reports both native and unified-input forecasting
comparisons, and makes the boundary between the two data sources explicit in
the interface.
